% =====================================================================
%  Carta formal — Projeto Hermathena / BINGO
%  Compilar com LuaLaTeX ou XeLaTeX (Overleaf: Menu > Compiler).
%  Mantenha hermathena.sty no mesmo diretório.
% =====================================================================
\documentclass[11pt,a4paper]{article}

\usepackage[utf8]{inputenc}
\usepackage[T1]{fontenc}
\usepackage[left=3cm,right=2.5cm,top=3cm,bottom=3cm]{geometry}
\usepackage{hermathena}
\usepackage{setspace}
\usepackage{parskip}

\onehalfspacing
\pagestyle{empty}
\setlength{\parindent}{0pt}

% ---------------------------------------------------------------------
%  Dados da carta — edite aqui
% ---------------------------------------------------------------------
\newcommand{\cartaRemetente}{Projeto Hermathena — Comunicação Multimídia do Radiotelescópio BINGO}
\newcommand{\cartaRemetenteDetalhe}{Universidade Federal de Campina Grande (UFCG)}
\newcommand{\cartaLocalData}{Campina Grande, 9 de outubro de 2026}
\newcommand{\cartaDestinatario}{À Direção do Instituto de Pesquisa Parceiro}
\newcommand{\cartaDestinatarioDetalhe}{Avenida Exemplo, 1000 — Cidade/UF}
\newcommand{\cartaAssunto}{Proposta de cooperação institucional}
\newcommand{\cartaSignatario}{Luciano Barosi}
\newcommand{\cartaSignatarioCargo}{Coordenador — Projeto Hermathena}

\begin{document}

% --- Cabeçalho -------------------------------------------------------
\begin{minipage}[t]{0.62\textwidth}
  \hermaTitle{\cartaRemetente}\\[2pt]
  \hermaSubtitle{\cartaRemetenteDetalhe}
\end{minipage}\hfill
\begin{minipage}[t]{0.33\textwidth}
  \raggedleft\hermathenaLogo[2.6cm]
\end{minipage}

\vspace{4pt}
\hermaRule
\vspace{18pt}

% --- Destinatário ----------------------------------------------------
\cartaDestinatario\\
\cartaDestinatarioDetalhe

\vspace{14pt}
\textbf{Assunto:} \cartaAssunto

\vspace{14pt}
Prezados(as) Senhores(as),

\vspace{6pt}
O Projeto Hermathena conduz a estratégia de comunicação, educação científica e
engajamento público associada ao Radiotelescópio BINGO, instalado em Aguiar, na
Paraíba. Suas ações articulam pesquisa de fronteira, formação de recursos humanos
e desenvolvimento regional, em diálogo permanente com a sociedade.

\vspace{6pt}
[Corpo da carta. Prefira a voz ativa, o presente do indicativo e um assunto por
parágrafo. Defina siglas na primeira ocorrência e referencie as fontes sempre que
houver afirmação factual.]

\vspace{6pt}
Colocamo-nos à disposição para detalhar esta proposta e para construir, em
conjunto, um desenho institucional de cooperação.

\vspace{24pt}
Respeitosamente,

\vspace{28pt}
\rule{7cm}{0.4pt}\\
\textbf{\cartaSignatario}\\
\cartaSignatarioCargo

\vspace{20pt}
\hermaRule
\vspace{4pt}
{\footnotesize\color{hermaGrafite} Documento institucional do Projeto Hermathena ·
Radiotelescópio BINGO · \cartaLocalData}

\end{document}
